\chapter{Conclusioni}
\label{sec:conclusions}

L'obiettivo di questa tesi è stato trasporre una selezione di programmi aggregati, originariamente sviluppati in Proto, nel linguaggio Collektive, con l'intento di offrire una guida pratica alla migrazione di sistemi legacy verso soluzioni più moderne e accessibili. 

Il percorso intrapreso per raggiungere questo obiettivo ha richiesto un'analisi approfondita della programmazione aggregata e della sua evoluzione storica. Partendo dai concetti fondamentali (\Cref{chap:introduzione}), sono state esaminate le basi teoriche di \ac{FC} e \ac{XC}, per poi tracciare l'evoluzione dei linguaggi che hanno implementato questi principi. Particolare attenzione è stata dedicata a Proto (\Cref{chap:mitproto}), linguaggio pionieristico che ha posto le fondamenta della programmazione aggregata, analizzandone sintassi e costrutti. L'analisi di Collektive (\Cref{chap:collektive}) ha poi rivelato come questo linguaggio non solo erediti i principi teorici del Field Calculus e di XC, ma li integri in un contesto moderno e multipiattaforma.

Attraverso la trasposizione di una serie di programmi rappresentativi (\Cref{chap:trasposizione}), si è dimostrato che Collektive è in grado di rappresentare efficacemente le logiche dei programmi originali, migliorandone la leggibilità e semplificandone lo sviluppo grazie a una libreria standard ricca e versatile. Tale libreria fornisce funzioni auto-stabilizzanti che consentono ai programmi di ereditare automaticamente proprietà di resilienza e adattabilità. Un ulteriore vantaggio significativo è la presenza nativa dell'operatore $\texttt{share}$, che permette di modellare elegantemente l'evoluzione spazio-temporale dei campi computazionali in un unico costrutto, ottenendo prestazioni superiori rispetto all'approccio di Proto che richiede la combinazione esplicita di $\texttt{rep}$ e $\texttt{nbr}$.


Questo lavoro apre molte possibilità per rafforzare ulteriormente il ruolo di Collektive nell'ecosistema della programmazione aggregata. Innanzitutto, un confronto sistematico tra Collektive e altri linguaggi come Protelis, ScaFi e FCPP fornirebbe agli sviluppatori una valutazione oggettiva di performance, espressività ed ergonomia. Trasponendo gli stessi esempi in tutti questi framework, sarebbe possibile offrire una guida preziosa per la scelta dello strumento più adatto. Inoltre, lo sviluppo di un "Collektive Playground", sul modelli di  WebProto e ScaFi Web, permetterebbe agli sviluppatori di sperimentare direttamente nel browser, accelerando l’adozione del framework e facilitandone l’apprendimento.

In conclusione, l'obiettivo è stato raggiunto, dimostrando che Collektive rappresenta un'evoluzione matura nel campo della programmazione aggregata. 
%
Fornendo una ``stele di Rosetta'' per la traduzione da Proto, non solo si valida l'espressività del nuovo linguaggio, ma si offre anche una risorsa concreta alla comunità, abbassando le barriere all'ingresso e promuovendo la fiducia verso questa tecnologia.
%
Il codice prodotto è stato integrato nel repository ufficiale degli esempi di Collektive, contribuendo direttamente all'ecosistema del progetto e garantendo che i risultati di questo lavoro possano beneficiare futuri sviluppatori.
